\chapter{Talking to a Database}
\label{ch:database}

\standfirst{SQL through ODBC, a query builder that finds its own joins, exact
decimals, and a query that does not stop the other cores while it runs.}

The \ct{Database} category reaches PostgreSQL, MySQL, SQLite, Oracle and SQL
Server. It is the first part of this system that talks to something outside
itself, and the only part that needs a library the Blue Book never heard of.

\section{A first query}

\begin{st}
DbConnection open: 'DRIVER=SQLITE3;Database=payroll.db;' do: [:db |
    db execute: 'INSERT INTO employee VALUES (?, ?, ?)'
       with: (Array with: 1 with: 'Ada' with: 'Lovelace').
    db commit.
    (db fetchOne: 'SELECT last_name FROM employee WHERE employee_id = ?'
        with: (Array with: 1)) at: 'last_name']
\end{st}

Three things in that are worth naming straight away.

The \ct{?} marks are \emph{bound parameters} and the values go beside them,
never into the text. The column is read by name, case-insensitively, because
SQL identifiers are. And \ct{open:do:} closes the connection even if the block
raises---which is the form to prefer, because a connection left open by a
method that returned early is a server-side resource nothing in this image
will ever free.

\section{One connection per process}

A \ct{DbConnection} belongs to one process. That is not a limitation of the
implementation; it is the design, and it is what makes database access here
genuinely parallel rather than merely concurrent.

Every other Smalltalk schedules its processes green, so $N$ processes share
one connection and take turns at one socket. Here $N$ workers hold $N$
connections on $N$ cores and the database sees $N$ clients.

\begin{stwarn}
Sharing one connection between two processes is a program error. A driver has
one result set per statement and no notion of whose turn it is, so the second
caller does not wait---it corrupts the first one's answer. Open one per
process. They are cheap next to what they do.
\end{stwarn}

\section{Transactions}

Autocommit is \emph{off}. A connection opens inside a transaction and nothing
is durable until you say so:

\begin{st}
db transact: [ record update. other delete ]
\end{st}

\ct{transact:} commits when the block finishes and rolls back if it raises.
Prefer it over \ct{commit} and \ct{rollback} by hand: an exception through a
half-finished transaction otherwise leaves the connection holding locks the
server will not release until somebody decides, and this decides.

Off by default is the safe direction. A program that forgets to commit loses
work it did not mean to keep; autocommit loses the ability to undo work it did
not mean to do.

\section{Reading}

\ct{query:} answers a \ct{DbCursor}, which reads one row at a time and closes
itself when it reaches the end:

\begin{st}
(db query: 'SELECT * FROM employee ORDER BY last_name')
    do: [:each | Transcript show: (each at: 'last_name'); cr]
\end{st}

\ct{do:}, \ct{collect:}, \ct{select:} and \ct{fetchAll} all close the cursor
when they finish---including when the block leaves early or raises, because
the close is in an \ct{ensure:}. A leaked statement is invisible until a
connection runs out of them.

\begin{center}
\small
\begin{tabular}{@{}ll@{}}
\toprule
\ctp{next} & the next row as a \ctp{DbRecord}, or nil\\
\ctp{fetchOne} \ctp{fetchAll} & one row, or all of them, then close\\
\ctp{do:} \ctp{collect:} \ctp{select:} \ctp{detect:ifNone:} \ctp{inject:into:} & and close\\
\ctp{columns} \ctp{columnCount} & the \ctp{DbColumnInfo} of each column\\
\ctp{maxRows:} & stop after this many\\
\ctp{close} \ctp{isClosed} & harmless twice\\
\bottomrule
\end{tabular}
\end{center}

\begin{stnote}[Every row is a new record]
The Java this was ported from reuses one row object across a cursor. Here the
process holding a row and the process advancing the cursor can be on two cores
at once, so reuse would turn \ct{cursor fetchAll} into a collection of $N$
references to one row holding the last one's values. An allocation per row is
small beside the round trip that produced it.
\end{stnote}

\subsection{What a column answers}

\begin{center}
\small
\begin{tabular}{@{}ll@{}}
\toprule
\ctp{TINYINT} \ldots{} \ctp{BIGINT} & \ctp{Integer}, exactly, at any width\\
\ctp{REAL} \ctp{FLOAT} \ctp{DOUBLE} & \ctp{Float}\\
\ctp{DECIMAL} \ctp{NUMERIC} & \ctp{Fraction} --- exact\\
\ctp{CHAR} \ctp{VARCHAR} and the long and national forms & \ctp{String}\\
\ctp{BINARY} \ctp{VARBINARY} \ctp{BLOB} & \ctp{ByteArray}\\
\ctp{DATE} & \ctp{Date}\\
\ctp{TIME} & \ctp{Time}\\
\ctp{TIMESTAMP} & \ctp{DbTimestamp}\\
\ctp{BIT} \ctp{BOOLEAN} & \ctp{true} / \ctp{false}\\
\ctp{NULL} & \ctp{nil}\\
\bottomrule
\end{tabular}
\end{center}

A record also has typed accessors for the cases where you know better than the
column does---\ct{integerAt:}, \ct{numberAt:}, \ct{stringAt:},
\ct{booleanAt:}, \ct{dateAt:}, \ct{bytesAt:} and \ct{dateIntegerAt:}, which
answers a date as \ctp{yyyymmdd}.

\ct{at:} raises if there is no such column, and that is deliberate: a
misspelled name answering nil is indistinguishable from a \ctp{NULL}, and the
report that comes out is wrong rather than missing. Use \ct{at:ifAbsent:} when
absence is a real possibility.

\ct{DbTimestamp} exists because the 1983 image has no single class for a date
and a time together, and its \ct{Time} is whole seconds. A timestamp is both
halves at once and carries a fraction. It answers \ct{asDate}, \ct{asTime},
\ct{asInteger}, and \ct{asDateAndTime} in an image that has Pharo's
\ct{DateAndTime} loaded.

\section{Writing}

A record from a single-table select knows where it came from and can be
written back:

\begin{st}
| row |
row := db fetchOne: 'SELECT * FROM employee WHERE employee_id = 1'.
row at: 'last_name' put: 'Byron'.
row update                            "answers the number of rows changed"
\end{st}

\begin{center}
\small
\begin{tabular}{@{}ll@{}}
\toprule
\ctp{insert} & \ctp{INSERT}, answering the count\\
\ctp{insertReturningKey:} & and answer the key the database generated\\
\ctp{update} & only the columns that changed\\
\ctp{delete} & matched on the primary key\\
\ctp{copyCorrespondingFrom:} & take every column both records have\\
\bottomrule
\end{tabular}
\end{center}

A new row is built with \ct{newRecord:} and filled in:

\begin{st}
| record |
record := db newRecord: 'employee'.
record at: 'first_name' put: 'Edsger'; at: 'last_name' put: 'Dijkstra'.
record insertReturningKey: 'employee_id'
\end{st}

\ct{insertReturningKey:} is where the five databases differ and the only place
in this package that has to know which one it is talking to. PostgreSQL gets
\ctp{RETURNING}, SQLite \ctp{last\_insert\_rowid()}, MySQL
\ctp{LAST\_INSERT\_ID()}, SQL Server \ctp{SCOPE\_IDENTITY()}. ODBC has no
equivalent of JDBC's \ctp{getGeneratedKeys} and this is the substitute.

\begin{stnote}[Update writes only what changed]
Writing every column---which is what the Java does---means two processes
updating two different fields of one row each overwrite the other's work, and
neither ever sees a conflict. In a green-threaded Smalltalk that race needs
two transactions to interleave; here it needs two cores.

The key is matched \emph{as it was read}, not as it now stands. A record whose
primary key is itself changed must still be found by its old key, or the
update matches nothing---or, worse, matches a different row that has since
taken the new value.
\end{stnote}

A record from a join has no single table to be written back to, and says so
rather than guessing. So does a table with no primary key: there is no way to
name one row of it, and an update built without a key would match every row in
the table.

\section{The query builder}

The reason the category is worth having. You say what you want; it works out
how to reach it:

\begin{st}
(db newQueryBuilder: 'employee')
    select: 'employee.last_name';
    select: 'department.name';
    where: 'project.project_id = ?' with: (Array with: 17);
    orderBy: 'employee.last_name';
    fetchAll
\end{st}

Nothing there says how \ctp{employee} reaches \ctp{department}, or how either
reaches \ctp{project}. The builder reads the table names out of the columns
and conditions, asks the schema graph for the foreign keys that connect them,
and writes the joins:

\begin{plain}
SELECT employee.last_name, department.name FROM employee
  JOIN department ON employee.department_id = department.department_id
  JOIN assignment ON employee.employee_id = assignment.employee_id
  JOIN project ON assignment.project_id = project.project_id
 WHERE project.project_id = ? ORDER BY employee.last_name
\end{plain}

\ctp{assignment} appears in that and nowhere in the query that asked for it.
Adding a column from a fourth table adds a join without anybody editing a
\ctp{FROM} clause, which is the point: a hand-written join that is one table
out of date is a query that still runs.

\begin{center}
\small
\begin{tabular}{@{}ll@{}}
\toprule
\ctp{select:} \ctp{selectAll:} \ctp{distinct} & columns\\
\ctp{selectCount:as:} \ctp{selectSum:as:} \ctp{selectAverage:as:} & aggregates\\
\ctp{selectMinimum:as:} \ctp{selectMaximum:as:} & \\
\midrule
\ctp{where:with:} & a condition and its values\\
\ctp{where:in:} \ctp{where:notIn:} & an \ctp{IN} list, bound\\
\ctp{whereExists:} \ctp{whereNotExists:} & taking another builder\\
\ctp{startOr} \ctp{startAnd} \ctp{endGroup} & nesting\\
\midrule
\ctp{groupBy:} \ctp{having:with:} \ctp{orderBy:} \ctp{orderByDescending:} & \\
\ctp{limit:} & spelled per database\\
\midrule
\ctp{leftJoin:} \ctp{rightJoin:} & reach a table this way instead\\
\ctp{join:column:to:column:} & a join the graph does not have\\
\ctp{table:as:} & an alias, for a self-join\\
\ctp{with:as:} \ctp{union:} \ctp{unionAll:} & \\
\midrule
\ctp{build} \ctp{parameters} & the SQL, and its values\\
\ctp{query} \ctp{fetchAll} \ctp{fetchOne} \ctp{do:} & run it\\
\bottomrule
\end{tabular}
\end{center}

\subsection{Grouping conditions}

\ctp{a AND (b OR c)} is not \ctp{a AND b OR c}, and a builder that kept
conditions in a flat list could only produce the second:

\begin{st}
(db newQueryBuilder: 'employee')
    where: 'employee.active = ?' with: (Array with: 1);
    startOr;
        where: 'employee.last_name = ?' with: (Array with: 'Hopper');
        where: 'employee.last_name = ?' with: (Array with: 'Lovelace');
    endGroup
\end{st}

\begin{plain}
SELECT * FROM employee WHERE employee.active = ?
  AND (employee.last_name = ? OR employee.last_name = ?)
\end{plain}

The values come out in the order the question marks are \emph{written}, not
the order the methods were called in---each condition carries its own values,
so moving a clause moves its parameters with it.

\subsection{Nothing is pasted into the text}

\begin{stwarn}[Injection]
Not one value written through this package is ever put into SQL text. Every
one goes as a bound parameter, so there is nothing for a quote in a surname to
break and nothing for an attacker to close early. The only things built into
the text are table and column names, and those come from the schema.

That is why every method taking a value takes it separately:

\begin{st}
q where: 'salary > ?' with: (Array with: amount)   "safe, with any value"
q where: 'salary > ', amount printString           "not, and never"
\end{st}
\end{stwarn}

An empty \ct{where:in:} becomes a condition that is false rather than
\ctp{IN ()}, which is a syntax error in every database---so an empty filter
produces an empty result instead of a broken query.

\section{The join graph}

\ct{DbSchemaGraph} models the schema as a graph: tables are nodes, foreign
keys are edges. A connection reads its own on first use and keeps it, because
walking a large schema's metadata is a great many round trips.

You can also build one by hand, which is what to do when the database was
never told about a relationship that exists:

\begin{st}
| graph |
graph := DbSchemaGraph new.
graph from: 'employee' column: 'department_id'
       to: 'department' column: 'department_id'.
graph from: 'order_line' columns: #('order_id' 'product_id')
       to: 'order_product' columns: #('order_id' 'product_id').
db schemaGraph: graph
\end{st}

Composite keys are the ordinary case rather than a special one, and the two
forms are the same method with one column or several.

Connecting $N$ tables with the fewest joins is the Steiner tree problem, which
is NP-hard, so this does what everyone does: take one table as the root, then
repeatedly attach whichever remaining table has the shortest breadth-first
path to anything already attached. That is an approximation and it is not
always the true minimum. It is \emph{deterministic}, which matters more here:
the same query must produce the same SQL every time, or two runs of one report
join differently and nobody can tell which was right.

Two tables with no path between them raise, rather than producing a cross join
nobody asked for and a result set the size of the product.

A graph can be written out and read back---\ct{storeOn:} and \ct{readFrom:},
one line per edge---so an application can read the schema once and keep the
answer.

\section{Exact decimals}

\ctp{DECIMAL} and \ctp{NUMERIC} columns come back as the digits the database
printed, read into a \ct{Fraction}:

\begin{st}
(Odbc numberFromDigits: '0.10') * 10        "1, exactly"
'0.10' asNumber * 10                        "1.0, and never was one tenth"
\end{st}

This is the one respect in which the port is better than the Java it came
from. A money column read through a float is wrong in exactly the cases
anybody cares about, and wrong quietly. Smalltalk has exact arithmetic, so the
only place the exactness could be lost is this boundary---and the boundary is
where it is easiest not to lose it.

The image's own \ct{Number class>>readFrom:} converts the fractional part with
\ct{asFloat} and divides. That is 1983's behaviour and it stays; this package
does not go through it.

\begin{stwarn}[SQLite has no decimal type]
SQLite applies numeric affinity: a column declared \ctp{DECIMAL(20,4)} stores
\ctp{12345678901234.5678} as a C double and hands back
\ctp{12345678901234.6}. The digits are gone at \ctp{INSERT}, before anything
here sees them, and its driver then reports the column as \ctp{VARCHAR}---so
this package correctly answers the string it was given, and there was never an
exact value to answer.

Checking the claim above on SQLite will suggest it does not hold. It is SQLite
that does not hold it. PostgreSQL, Oracle, SQL Server and MySQL all have a
real \ctp{DECIMAL} and report it as one.
\end{stwarn}

\section{It does not stop the world}

This is the part that is specific to this system, and the reason the category
could not simply be lifted from somewhere else.

A worker inside \ctp{SQLExecute} is not running bytecodes, so it never reaches
the safepoint poll (Chapter~\ref{ch:contract}). Nothing is wrong with that
until a collection is asked for, and then everything is: the collector waits
for a worker that is waiting for a server, and one slow query stops every core
for as long as the query takes.

So every database call that can block declares itself. The worker parks for
the duration exactly as the poll would---counted among the parked, registers
written back---and the collector runs \emph{while} it sits in the driver.

Measured, by deleting the declaration and running the same test:

\begin{center}
\small
\begin{tabular}{@{}lll@{}}
\toprule
& query & safepoint granted in\\
\midrule
with the native regions & 0.24s & \textbf{0.0000s}\\
without them & 0.24s & \textbf{0.2396s}\\
\bottomrule
\end{tabular}
\end{center}

Without them the collector waits for the whole query. With them it does not
wait at all. \ct{tests/unit/test_odbc_parallel.c} is that measurement, and
it fails if the property regresses.

\begin{stnote}
Nothing you write has to do anything about this. It is a property of the
primitive, not a rule for callers. It is here because it is the answer to
``what happens to the other cores while my report runs'', and the answer is
``they keep running''.
\end{stnote}

\section{Getting a driver}

ODBC is optional and its absence is a first-class outcome. A build without it
has every method present, \ct{Odbc isAvailable} answering false, and any
attempt to connect raising with a sentence saying so.

\begin{sh}
make deps                # says whether ODBC was found, and how to install it
make NODB=1              # build without it, on a machine that has it
\end{sh}

Fedora wants \ct{unixODBC-devel}, Debian \ct{unixodbc-dev}; macOS ships iODBC
and Homebrew has unixODBC, and both are found. Windows carries the driver
manager in the operating system.

The connection string is ODBC's, not this system's---a DSN name, or the
driver-specific form for a driver that needs no DSN:

\begin{plain}
DSN=payroll;UID=blake;PWD=secret
DRIVER=SQLITE3;Database=payroll.db;
DRIVER={PostgreSQL Unicode};Server=localhost;Database=payroll;UID=blake;
\end{plain}

\ct{odbcinst -q -d} lists the drivers a machine has.

\section{The classes}

\begin{center}
\small
\begin{tabular}{@{}ll@{}}
\toprule
\ctp{DbConnection} & one per process; transactions, catalogue, caches\\
\ctp{DbCommand} & a prepared statement, reused across rows\\
\ctp{DbCursor} & a result set, one row at a time\\
\ctp{DbRecord} & one row, and the ability to write it back\\
\ctp{DbQueryBuilder} & SQL from what you want, not how to reach it\\
\ctp{DbSchemaGraph} & tables as nodes, foreign keys as edges\\
\ctp{DbForeignKey} & one key, composite or not\\
\ctp{DbColumnInfo} & what a result set says about a column\\
\ctp{DbTimestamp} & a moment, to the precision a database keeps one\\
\ctp{DbWhereGroup} & a node of a \ctp{WHERE} tree\\
\ctp{DbError} & what the driver said, in the driver's own words\\
\ctp{Odbc} & the only class that knows a primitive exists\\
\bottomrule
\end{tabular}
\end{center}

Everything above \ct{Odbc} is ordinary Smalltalk sending ordinary messages.
The category can be read without knowing ODBC, a second driver could be put
underneath without touching anything above it, and there is exactly one file
to look in when a database does something a database should not.

A \ct{DbError} carries the driver's own text, \ctp{SQLSTATE} and all. There is
no hierarchy of database error classes, because the useful distinction is not
between kinds of refusal---a caller almost never wants to handle a constraint
violation and not a deadlock---but between what the driver said and anything
this image might invent instead. If you do want to tell one refusal from
another, the \ctp{SQLSTATE} at the front of the message is the part that is
the same across every database.

\begin{sttry}
No database to hand? The query builder and the join graph need none---they
build strings. Build a graph by hand, point a builder at it with
\ct{DbQueryBuilder onGraph: g table: 'employee'}, and \ct{build} the SQL. That
is how the 43 tests in \ct{lib/Database-Tests} check the generator, and they
run everywhere.
\end{sttry}

\section{Where the rest of it is}

\ct{doc/DATABASE.md} carries the design notes: why ODBC rather than a driver
per database, the two things ODBC does not have, what the tests found, and the
provenance. \ct{lib/Database-Live-Tests} is the suite that needs a real
database, in a profile of its own, run deliberately:

\begin{sh}
./st80 -bootstrap -profile profiles/database-live.profile -tests
\end{sh}
